\documentclass[10pt,openany]{book}
\pagestyle{plain}

\usepackage{amsmath, amssymb, amsfonts, amsthm, mathrsfs, mathtools, mathalpha}
\usepackage{bbold}
\DeclareSymbolFont{AMSb}{U}{msb}{m}{n}
\DeclareSymbolFontAlphabet{\mathbb}{AMSb}
\DeclareMathAlphabet{\bboldmath}{U}{bbold}{m}{n}
\newcommand{\dsone}{\bboldmath{1}}
\usepackage{thmtools}
\usepackage{enumitem}
\usepackage[colorlinks=true, linkcolor=red, citecolor=blue]{hyperref}
\usepackage{cleveref}
\usepackage{diagbox, array}
\usepackage{graphicx}
\usepackage{xparse}
\usepackage{url}
\usepackage{float}

% \usepackage{fontspec}
% \newfontfamily\devanagarifont[Script=Devanagari]{Sanskrit 2003}

\usepackage{subcaption}
\usepackage[justification=centering]{caption}

\usepackage{tikz}
\usetikzlibrary{graphs, automata, positioning}

\newtheorem{theorem}{Theorem}[chapter]
\newtheorem{proposition}[theorem]{Proposition}
\newtheorem{conjecture}[theorem]{Conjecture}
\newtheorem{lemma}[theorem]{Lemma}
\newtheoremstyle{italicnonumber}{}{}{\itshape}{}{\itshape}{}{.5em}{}
\theoremstyle{italicnonumber}
\newtheorem*{slemma}{Lemma}
\newenvironment{prooflem}[1][]{
  \renewcommand{\qedsymbol}{\ensuremath{\backslash\backslash\backslash}}
  \begin{proof}[Proof of Lemma\ifx\\#1\\\else~#1\fi]
}{
  \end{proof}
  \renewcommand{\qedsymbol}{\oldqedsymbol}%
}
\newtheorem{corollary}[theorem]{Corollary}
\newtheorem{exercise}[theorem]{Exercise}
\theoremstyle{definition}
\newtheorem{definition}[theorem]{Definition}
\newtheorem{notation}[theorem]{Notation}
\newtheorem{example}[theorem]{Example}
\newenvironment{question}{\par\vspace{1ex}\noindent\textbf{Question.}\ }{\par}
\newenvironment{answer}{\par\noindent\emph{Ans.}\ }{\hfill$\square$\par\vspace{1ex}}
\usepackage [english]{babel}
\usepackage [autostyle, english = american]{csquotes}
\MakeOuterQuote{"}
\renewcommand\qedsymbol{$\blacksquare$}
\let\oldforall\forall
\renewcommand{\forall}{\;\oldforall\;}
\let\oldexists\exists
\renewcommand{\exists}{\;\oldexists\;}

\newcommand{\C}[2]{\binom{#1}{#2}}
\newcommand{\M}[2]{\left(\!\! \begin{array}{c} #1 \\ #2 \end{array} \!\!\right)}
\newcommand{\cb}[2]{\genfrac{\{}{\}}{0pt}{}{#1}{#2}}
\newcommand{\sqb}[2]{\genfrac{[}{]}{0pt}{}{#1}{#2}}

\newcommand{\N}{\mathbb{N}}
\newcommand{\Z}{\mathbb{Z}}
\newcommand{\Q}{\mathbb{Q}}
\newcommand{\R}{\mathbb{R}}
\newcommand{\bC}{\mathbb{C}}
\newcommand{\E}{\mathbb{E}}
\newcommand{\bP}{\mathbf{P}}
\newcommand{\mc}[1]{\mathcal{#1}}
\newcommand{\ms}[1]{\mathscr{#1}}
\newcommand{\emp}{\emptyset}
\newcommand{\ve}{\varepsilon}
\newcommand{\su}{\subseteq}
\newcommand{\pr}{\prime}
\newcommand{\io}{\iota}
\newcommand{\lan}{\langle}
\newcommand{\ran}{\rangle}
\newcommand{\la}{\lambda}
\newcommand{\om}{\Omega}
\newcommand{\fk}{\mathfrak}

\usepackage{geometry}
\geometry{margin=1in}
\setlength{\parskip}{0.5em}
\setlength{\parindent}{0pt}

\raggedbottom

\usepackage[backend=biber]{biblatex}
\addbibresource{references.bib}

\begin{document}

\title{MA523: Basic Number Theory notes}

\author{\emph{Om Joglekar}}
\date{}
\maketitle

\tableofcontents

\chapter{Lecture 1: 29th July, 2025}

We begin with the fundamental theorem of arithmetic.

\begin{theorem} [The Fundamental Theorem of Arithmetic] \label{funda_arithmetic}
  Every integer $n \geq 2$ can be written uniquely (up to order) as $n=p_1^{\alpha_1} p_2^{\alpha_2} \cdots p_k^{\alpha_k}$ for prime numbers $p_i$ and some powers $\alpha_i \geq 0$.
\end{theorem}

\begin{proof}
  One of the proofs is directly using concepts from basic ring theory. We consider the ring of integers with usual addition and multiplication. It is well established that this ring is a Euclidean domain. Thus, it is a PID and hence also a UFD.

  Let $I$ be an ideal. Since $\Z$ is a PID, $I=(n)$ for some $n \in \Z$. If we demand that $I$ is a prime ideal, it is clear that $n$ must not have any non-trivial divisors. Thus, prime ideals correspond to prime numbers in this ring.

  Now in a UFD, we know that an element is irreducible if and only if it is prime. By the very definition of a UFD, every non-zero non-unit can be written as a product of irreducible elements and this decomposition is unique up to order and scaling by units.

  Thus, every $n \geq 2$ can be expressed uniquely as a product of primes up to order. The only units in the ring are $1$ and $-1$ and $(-1)^2=1$ will imply that for $n \geq 2$, we will not come across negative integers in the decomposition.
\end{proof}

We now discuss some stuff related to prime numbers. There are some classical results and some important conjectures in the theory that follows.

\begin{proposition}
  There are infinitely many prime numbers.
\end{proposition}

\begin{proof}
  Suppose not. Let there be finitely many prime numbers $p_1, p_2, \dots, p_k$. Consider the number $p = p_1p_2\cdots p_k + 1$. It is clearly not divisible by any of the finitely many prime numbers we had. Thus, the only divisors of $p$ are $p$ and $1$ making it prime. This contradicts our assumption.
\end{proof}

\begin{proposition}
  There are infinitely many primes of the form $4k+3$.
\end{proposition}

\begin{proof}
  Suppose not. Let $p_1,\dots,p_k$ be finitely many primes of the form $4k+3$. Consider the number $N = 4 \left( p_1 \cdots p_k \right)-1$.

  It is clear that $N = 3 \; (\text{mod } 4)$. Observe that the product of two integers of type $4k+1$ is again an integer of type $4k+1$. Since $N$ is odd, it only has odd primes in its unique factorisation. By the observation we just made, it is not possible that all of these prime factors are of the form $4k+1$. Thus, there is at least one prime factor of the form $4k+3$. But this factor cannot be one of $p_1,p_2,\dots,p_k$ since $N$ is clearly not divisible by any of these. Thus, the only possible case it that $N$ is prime itself but this is a contradiction to the finite list of primes of the form $4k+3$ we had assumed.
\end{proof}

One may want to attempt a $4k+1$ proof but this is much more difficult compared to the one above.

A very powerful theorem was given by Dirichlet in 1837. The proof is very hard and lengthy and we will simply state the theorem.

\begin{theorem} [Dirichlet's Theorem on Arithmetic Progressions]
  Let $a$ and $b$ be coprime positive integers. There are infinitely many primes of the form $a+nb$.
\end{theorem}

We now have two famous theorems - Euler's theorem and Wilson's theorem

\begin{theorem} [Euler's Theorem] \label{eul_theo}
  Let $n$ be a positive integer and $a \in \Z$ such that $\gcd(a,n)=1$. Then $a^{\varphi(n)} = 1 \; (\text{mod } n)$ where $\varphi(n)$ denotes the Euler's totient function defined later in \cref{eu_tot}.
\end{theorem}

\begin{proof}
  Consider the group $U$ of units of $\Z/n\Z$ under multiplication. It is clear that $x \in \Z/n\Z$ is a unit if and only if $\gcd(x,n)=1$ and hence $\#U = \varphi(n)$. 
  
  By Lagrange's theorem, the order of every element divides the order of the group. Since $a \in U$, we have $a^{\varphi(n)}=1$ and hence the result follows.
\end{proof}

We have an immediate corollary below.

\begin{corollary} [Fermat's Little Theorem]
  Let $p$ be a prime number and $a$ be a positive integer which is not a multiple of $p$. Then, $a^{p-1} = 1 (\text{mod } p)$
\end{corollary}

\begin{theorem} [Wilson's theorem]
  For any prime number $p$, $(p-1)!+1$ is divisible by $p$.
\end{theorem}

\begin{proof}
  Consider the multiplicative group $\left( \Z/p\Z\right)^{\times}$. Every element $x$ has an inverse $x^{-1}$ in the group. The only elements which are self-inverses are the solutions of $x^2 = 1 (\text{mod } p)$. One can check that the only solutions are $x = \pm 1 = 1, p-1$.

  Thus, in the product $(p-1)!= 1 \cdot 2 \cdot \cdots \cdot (p-1)$ every term pairs with its inverse to produce $1$, modulo $p$. Thus, when viewed modulo $p$, $(p-1)!+1 = (1 \cdot (p-1)) + 1 = p$ and we are done. 
\end{proof}

We now have the famous Euler product for the Riemann-Zeta function.

\begin{proposition} \label{eul_prod}
  For some real number $s > 1$,
  \[
  \zeta(s) = \sum_{n=1}^\infty \frac{1}{n^s} = \prod_{p \text{ prime}} \left( 1-\frac{1}{p^s} \right)^{-1}
  \]
\end{proposition}

\begin{proof}
  The series $\zeta(s)$ clearly converges absolutely. We can thus open the RHS completely, rearrange and prove the proposition.

  Using series expansion of $\left( 1-r \right)^{-1}$ for $ |r| <1$, we get:
  \[
  \begin{aligned}
    \prod_{p \text{ prime}} \left( 1-\frac{1}{p^s} \right)^{-1} & = \prod_{p \text{ prime}} \left( \sum_{k=0}^\infty \frac{1}{p^{ks}} \right) \\
    & = \left( 1+\frac{1}{2^s} + \frac{1}{4^s} + \frac{1}{8^s} + \cdots \right) \left( 1+\frac{1}{3^s} + \frac{1}{9^s} + \frac{1}{27^s} + \cdots \right) \cdots
  \end{aligned}
  \]
  When multiplied out, each term is of the form $1/n^s$ and all possible $n \geq 1$ are covered by \cref{funda_arithmetic} proving the claim.
\end{proof}

Before moving on to the conjectures, we state Bertrand's postulate without proof. A proof may be found on the wikipedia page.

\begin{proposition} [Bertrand's Postulate]
  For any given integer $n \geq 2$, there exists a prime number in the interval $(n,2n)$.
\end{proposition}

We now state some famous conjectures on which lot of work has been done but they have not been proved - yet!

\begin{conjecture} [Goldbach]
  Every integer $ n > 2$ is the sum of two primes.
\end{conjecture}

\begin{conjecture} [Twin Prime]
  There are infinitely many pairs of primes $(p,p+2)$.
\end{conjecture}

\begin{conjecture} [Legendre]
  There is at least one prime in the interval $(n^2,(n+1)^2)$ for every $n \geq 1$.
\end{conjecture}

\begin{conjecture} [Polignac]
  For every even number $2k$, there are infinitely many prime pairs $(p,p+2k)$.
\end{conjecture}

\begin{conjecture} [Bunyakovsky]
  If a polynomial $f(n) \in \Z[n]$ is irreducible, has positive leading coefficient and does not have a fixed divisor $($some fixed integer $k > 1$ dividing all $f(n))$ then $f(n)$ takes infinitely many prime values.
\end{conjecture}

\chapter{Lecture 2: 1st August, 2025}

We begin with some notational clarifications through definition of asymptotic notations.

\begin{definition}
Let $f, g : \mathbb{R} \to \mathbb{R}$. We write $f(x) = O(g(x))$ on some set $S \su \R$ if there exist constants $C > 0$ such that
\[
|f(x)| \leq C |g(x)| \quad \forall x \in S
\]
\end{definition}

One sometimes writes $f(x) \ll g(x)$ to mean $f(x) = O(g(x))$.

\begin{definition}
Let $f, g : \mathbb{R} \to \mathbb{R}$ be functions. We write $f(x) \asymp g(x)$ on some set $S \su \R$ if $f(x) \ll g(x)$ and $g(x) \ll f(x)$, both hold on $S$.
\end{definition}

\begin{definition}
Let $f, g : \mathbb{R} \to \mathbb{R}$ and $x_0 \in \mathbb{R}$. We write $f(x) = o(g(x))$ as $x \to x_0$ if
\[
\lim_{x \to x_0} \frac{f(x)}{g(x)} = 0.
\]
\end{definition}

\begin{definition}
Let $f, g : \mathbb{R} \to \mathbb{R}$ and $x_0 \in \mathbb{R}$. We write $f(x) \sim g(x)$ as $x \to x_0$ if
\[
\lim_{x \to x_0} \frac{f(x)}{g(x)} = 1.
\]
\end{definition}

\begin{definition}
Let $f, g : \mathbb{R} \to \mathbb{R}$ and $x_0 \in \mathbb{R}$. We write $f(x) = \Omega(g(x))$ as $x \to x_0$ if
\[
\limsup_{x \to x_0} \left|\frac{f(x)}{g(x)}\right| > 0.
\]
\end{definition}

In the first two definitions, if the set $S$ is not mentioned (which is the most common case), we assume $S$ to be $[x_0,\infty)$ for some $x_0$ large enough chosen at our own discretion. In the last three definitions, if the point $x_0$ is not mentioned (which is also the most common case), we take it be $\infty$. Most often, we skip the set $S$ and the point $x_0$ because we only care about the limiting behaviour at infinity.

Below are some exercises to solidify the concepts of the above notations.

\begin{question}
Determine whether each of the following statements is true or false. Justify your answer.
\begin{enumerate}[label=(\roman*)]
    \item $x^{1000} = O(e^x)$
    \item $(\log x)^{1000} = O\left(\frac{1}{x^{100000}}\right)$
    \item $x^5 + 3x^2 + 7 \asymp x^5$
    \item $\dfrac{1}{x} \asymp \dfrac{1}{x^2}$
    \item $\log x = o(x^\varepsilon)$ for every $\varepsilon > 0$
    \item $x^2 + \log x = o(x^2)$
    \item $x^2 \sim x^2 + (\log x)^{1000}$
    \item $\dfrac{x}{\log x} \sim x$
     \item $x^3 + 5x = \Omega(x^3)$
    \item $\dfrac{1}{x} = \Omega\left(\dfrac{1}{x^2}\right)$
\end{enumerate}
\end{question}

\begin{answer}
\begin{enumerate}[label=(\roman*)]
  \item True. Choose $x_0 = 1000$ and $C=(1000)!$. The inequality $x^{1000} \leq C e^x$ is clearly true for all $x \geq x_0$.
  \item False. Suppose it were true, then there would exist a $C>0$ and an $x_0$ such that
  \[
  x^{100000} (\log{x})^{1000} \leq C \quad \forall x \geq x_0
  \]
  But as $x \to \infty$, the LHS also grows without bound and no such $C$ can exist.
  \item True. The following clearly holds for all $x \geq 4$
  \[
  |x|^5 \leq |x^5+3x^2+7| \leq |x|^5 + 3|x|^2 + 7 \leq 11 |x|^5
  \]
  \item False. Suppose some $\alpha, \beta >0$ exist and $x_0$ exists such that the following holds for all $x \geq x_0$.
  \[
  \alpha \frac{1}{x^2} \leq \frac{1}{|x|} \leq \beta \frac{1}{x^2}
  \]
  We get a contradiction in the right inequality by choosing $x > \max\{\beta + 1, x_0\}$.
  \item True. The limiting value is indeed $0$.
  \item False. The limiting value is $1$.
  \item True. The limiting value is $1$.
  \item False. The limiting value is $0$.
  \item True. The limiting value is $1$.
  \item True. The limiting value is $\infty$.
\end{enumerate}
\end{answer}

\begin{proposition}
  If $f(x) = o(g(x))$, then $f(x)=O(g(x))$.
\end{proposition}

\begin{proof}
  We know that the following limit holds
  \[
  \lim_{x \to \infty} \frac{f(x)}{g(x)} = 0
  \]
  Thus we have
  \[
  \lim_{x \to \infty} \frac{|f(x)|}{|g(x)|} = 0
  \]
  Pick any $C>0$. By definition of limit, we have the existence of an $x_0$ so that for all $x \geq x_0$ we have $|f(x)| \leq C |g(x)|$. Thus, $f(x) = O(g(x))$.
\end{proof}

\begin{question}
  Is the converse of the above statement true?
\end{question}

\begin{answer}
  No. It is false. Consider the functions $f(x)=g(x)=x$. Clearly, $f(x) = O(g(x))$ and $g(x) = O(f(x))$ but $f(x) \neq o(g(x))$ and $g(x) \neq o(f(x))$.
\end{answer}

We now define one of the most important concepts of analytic number theory - an arithmetical function

\begin{definition}
  Any function $f: \N \to \mathbb{C}$ is called an \textit{arithmetical function}. 
\end{definition}

One may think of this as a sequence. Note that we treat $0$ as an 'un-natural number' here, that is, $0 \notin \N$. This is however not the case in other fields of mathematics including combinatorics.

One of the premise of number theory is to analyse the behaviour of the partial sum function of different arithmetical functions.

As the first example, we discuss the prime-counting function.

\begin{definition}
  The \textit{characteristic function of primes} $f:\N \to \bC$ is defined as
  \[
  f(x) = \begin{cases}
    1 & \text{if } x \text{ is prime} \\
    0 & \text{otherwise}
  \end{cases}
  \]
\end{definition}

The partial sum function of the characteristic function of primes, also known as the \textit{prime-counting function}, $\pi:\N \to \bC$ given by $\pi(X) = \sum_{i=1}^X f(X)$ is the number of primes less than or equal to $X$. This is a famous function which has no closed form solution.

We now discuss many other examples of important arithmetical functions.

Given a number $n\in \N$, we can decompose it into its prime factorisation as $n=p_1^{\alpha_1} p_2^{\alpha_2} \cdots p_k^{\alpha_k}$ where $p_1,p_2,\dots,p_k$ are primes and $\alpha_1,\alpha_2,\dots,\alpha_k \geq 1$ are integers.

\begin{definition}
  The \textit{Möbius function} $\mu:\N \to \bC$ is defined as follows:
  \[
  \mu(n) = \begin{cases}
    (-1)^k & \text{if } n=p_1 p_2 \cdots p_k \geq 2 \\
    0 & \text{if } n \text{ is divisible by a square}, n \geq 2 \\
    1 & \text{if } n=1
  \end{cases}
  \]
\end{definition}

It is a well known conjecture that the partial sum function of the Möbius function is $O(X^{0.5+\varepsilon})$ for $\varepsilon > 0$ small enough.

\begin{definition} \label{eu_tot}
  The \textit{Euler's totient function} $\varphi : \mathbb{N} \to \bC$ is defined as
  \[
    \varphi(n) = \#\{k \in \mathbb{Z} \mid 1 \leq k \leq n, \, \gcd(k, n) = 1\}
  \]
\end{definition}

This function essentially returns the number of positive integers less than or equal to $n$ which are coprime to $n$. There is also a direct formula which we shall see soon.

\begin{definition}
  The \textit{von Mangoldt function} $\Lambda : \mathbb{N} \to \bC$ is defined as
  \[
    \Lambda(n) = 
    \begin{cases}
      \log p & \text{if } n = p^k \text{ for some prime } p \text{ and integer } k \ge 1 \\
      0 & \text{otherwise}
    \end{cases}
  \]
\end{definition}

\begin{definition}
  The \textit{divisor function} \( d : \mathbb{N} \to \bC \) is defined as
  \[
    d(n) = \#\{k \in \mathbb{N} \mid k \mid n\},
  \]
  i.e., \( d(n) \) is the number of positive integers dividing \( n \).
\end{definition}

\begin{definition}
  The \textit{Liouville function} \( \lambda : \mathbb{N} \to \bC \) is defined as
  \[
    \lambda(n) = (-1)^{\alpha_1 + \cdots + \alpha_k},
  \]
  where $n = p_1^{\alpha_1} p_2^{\alpha_2} \cdots p_k^{\alpha_k}$
\end{definition}

We now prove one of the central theorem of number theory (and combinatorics) - the Möbius inversion formula. We first need a lemma about the partial sums of the Möbius function.

\begin{lemma} \label{mob_lemma}
  We have,
  \[
  \sum_{d \mid n} \mu(d) = 0 \quad \forall n \geq 2
  \]
\end{lemma}

\begin{proof}
  Let $n=p_1^{\alpha_1} p_2^{\alpha_2} \cdots p_k^{\alpha_k}$. Since the Möbius function only takes non zero values on square free numbers, we only care about divisors $d$ of the number $p_1p_2 \cdots p_k$. To pick such a divisor, we need to pick a subset of these $p_i$ and the Möbius function will return the signed cardinality of the subset, that is, $-1$ is odd and $1$ if even.

  Since the number of odd size and even sized subsets of any finite set is equal, all signed cardinalities cancel out leaving us with 0.
\end{proof}

\begin{theorem}[Möbius Inversion Formula] \label{mob_inv}
  Let \( f, g : \mathbb{N} \to \mathbb{C} \) be arithmetic functions. Then
  \[
    g(n) = \sum_{d \mid n} f(d) \quad \text{for all } n \in \mathbb{N}
  \]
  if and only if
  \[
    f(n) = \sum_{d \mid n} \mu(d) \, g\left( \frac{n}{d} \right) \quad \text{for all } n \in \mathbb{N}.
  \]
\end{theorem}

\begin{proof}
Assume $ \displaystyle g(n)=\sum_{d\mid n}f(d)$ . Then 
\[
\sum_{d \mid n} \mu(d) g\left( \frac{n}{d}\right) = \sum_{d \mid n} \mu\left( \frac{n}{d}\right) g(d) = \sum_{d \mid n} \mu\left( \frac{n}{d}\right) \sum_{d' \mid d} f(d') = \sum_{d' \mid d \mid n} f(d') \mu\left( \frac{n}{d}\right)
\]

Now observe that if $d_1, d_2$ are divisors of $n$, then the following holds: $d_2 \mid d_1$ if and only if $\frac{n}{d_1} \mid \frac{n}{d_2}$ and if $m$ is some integer, $m \mid \frac{n}{d_2}$ if and only if $d_2 \mid \frac{n}{m} \mid n$.

Set $m=\frac{n}{d}$ and sum over $d'$ to get
\[
\sum_{d' \mid d \mid n} f(d') \mu\left( \frac{n}{d}\right) = \sum_{d' \mid n} f(d') \sum_{m \mid \frac{n}{d'}} \mu(m)
\]

By \cref{mob_lemma}, the second sum is $0$ everywhere except when $\frac{n}{d'} = 1$, that is, $n=d'$. Thus, the sum evaluates to $f(n)$.

The other direction is proved very similarly.
\end{proof}

We now obtain a formula for the Euler's totient function.

\begin{proposition}
  If $\varphi$ denotes the Euler's totient function, then for $n=p_1^{\alpha_1} p_2^{\alpha_2} \cdots p_k^{\alpha_k}$ we have,
  \[
  \varphi(n) = n\left(1-\frac{1}{p_1}\right)\left(1-\frac{1}{p_2}\right)\cdots \left(1-\frac{1}{p_k}\right)
  \]
\end{proposition}

\begin{proof}
  Consider, for $n$, the set $A(d) = \{1 \leq k \leq n : \gcd(k,n) = d\}$. One may check that as $d$ varies over the divisors of $n$, all the $A(d)$'s are mutually disjoint and their union is the entire set $[n]=\{ 1,2,\dots,n \}$.

  Further, the size of each $A(d)$ can be verified to be $\varphi\left(\frac{n}{d}\right)$ since every $k$ in $A(d)$ is of the form $k=d \ell$ for some $\ell \in [0, \frac{n}{d}]$ coprime to $\frac{n}{d}$. Conversely every such $k$ is present in $A(d)$ and hence the cardinality of $A(d)$ is the number of such $\ell$ which is precisely $\varphi\left(\frac{n}{d}\right)$.

  We thus have the identity
  \[
  n = \sum_{d \mid n} \varphi\left(\frac{n}{d}\right) = \sum_{d \mid n} \varphi(d)
  \]

  Applying \cref{mob_inv} with $g(n) = n$ and $f(n)=\varphi(n)$ we get
  \[
  \varphi(n) = \sum_{d \mid n} \mu(d) \frac{n}{d}
  \]

  It now suffices to prove the following identity:
  \[
  \sum_{d \mid n} \frac{\mu(d)}{d} = \left(1-\frac{1}{p_1}\right) \cdots \left(1-\frac{1}{p_k}\right)
  \]

  Since $\mu(d)$ is non-zero if and only if $d$ is square free, the summation reduces to the following
  \[
  \sum_{d \mid p_1p_2\cdots p_k} \frac{\mu(d)}{d} = \sum_{x_1=0}^1 \sum_{x_2=0}^1 \cdots \sum_{x_k=0}^1 \frac{\mu\left( p_1^{x_1} p_2^{x_2} \cdots p_k^{x_k} \right)}{p_1^{x_1} p_2^{x_2} \cdots p_k^{x_k}}
  \]
  We observe that $\mu(mn)= \mu(m)\mu(n)$ whenever $m$ and $n$ are coprime. This is easy to prove. Using this, we get
  \[
  \sum_{x_1=0}^1 \frac{\mu\left( p_1^{x_1} \right)}{p_1^{x_1}}\sum_{x_2=0}^1 \frac{\mu\left( p_2^{x_2} \right)}{p_2^{x_2}} \cdots \sum_{x_k=0}^1 \frac{\mu\left( p_k^{x_k} \right)}{p_k^{x_k}}
  \]

  Each $i$-th sum is individually equal to $ \displaystyle 1 - \frac{1}{p_i} $ and we are done.
\end{proof}

We used an important property of the Möbius function in the previous proof which motivates us to have the following definition.

\begin{definition}
  Let $f$ be a non-zero arithmetical function. We say $f$ is \textit{multiplicative} if $f(mn)=f(m)f(n)$ for coprime $m$ and $n$. We say that $f$ is \textit{completely multiplicative} if $f(mn)=f(m)f(n)$ for any $m$ and $n$.
\end{definition}

\begin{proposition}
  If $f$ is a multiplicative arithmetical function then we have
  \[
  \sum_{d \mid n} f(d) = \prod_{p \mid n, p \text{ prime}} \left[ 1 + f(p) + f(p^2) + \cdots + f(p^{r_p}) \right]
  \]
  where $r_p$ is the largest non-negative integer so that $p^{r_p} \mid n$.
\end{proposition}

\begin{proof}
  Let $n=p_1^{\alpha_1} p_2^{\alpha_2} \cdots p_k^{\alpha_k}$. We have
  \[
  \begin{aligned}
  \sum_{d \mid n} f(d)
  &= \sum_{x_1=0}^{\alpha_1} \sum_{x_2=0}^{\alpha_2} \cdots \sum_{x_k=0}^{\alpha_k}
  f(p_1^{x_1} p_2^{x_2} \cdots p_k^{x_k}) \\
  & = \sum_{x_1=0}^{\alpha_1} f(p_1^{x_1}) \sum_{x_2=0}^{\alpha_2} f(p_2^{x_2}) \cdots \sum_{x_k=0}^{\alpha_k} f(p_k^{x_k})
  \end{aligned}
  \]
  This is indeed the RHS, since $f(1)=1$ for multiplicative functions, and we are done.
\end{proof}

\begin{definition}
  Given two arithmetical functions $f$ and $g$, we define their \textit{Dirichlet convolution} to be the arithmetical function $f \star g$ given by $(f \star g)(n) = \sum_{d\mid n} f(d) g\left(\frac{n}{d}\right)$
\end{definition}

Under this definition, \cref{mob_inv} takes the following form: $g = f \star 1$ if and only if $f = \mu \star g$. Here, $1$ denotes the constant arithmetical function $1(n) = 1 \; \; \forall n$.

\begin{proposition}
  The set of all arithmetical functions $f$ such that $f(1) \neq 0$, forms an abelian group under the convolution operation.
\end{proposition}

\begin{proof}
  Firstly, it is very easy to check that the convolution operation is a binary operation on the set of multiplicative functions. $(f\star g)(1) = f(1)g(1) \neq 0$.

  The convolution operation is clearly commutative since $\sum_{d\mid n} f(d) g\left(\frac{n}{d}\right) = \sum_{d\mid n} f\left(\frac{n}{d}\right) g(d)$ clearly holds.

  We now verify the associativity:
  \[
  \begin{aligned}
  ((f\star g)\star h)(n) & = \sum_{d \mid n} (f \star g)(d) h\left(\frac{n}{d}\right) \\
  & = \sum_{d \mid n} \left( \sum_{d' \mid d} f(d') \; g\left( \frac{d}{d'} \right) \right) h\left(\frac{n}{d}\right) \\
  & = \sum_{d' \mid d \mid n} f(d') \; g\left( \frac{d}{d'} \right) \; h\left(\frac{n}{d}\right) \\
  & = \sum_{abc=n} f(a)g(b)h(c)
  \end{aligned}
  \]
  \[
  \begin{aligned}
  (f\star (g\star h))(n) & = \sum_{d \mid n} f\left( \frac{n}{d} \right) (g\star h)(d)  \\
  & = \sum_{d \mid n} f\left( \frac{n}{d} \right) \left( \sum_{d' \mid d} h(d') \; g\left( \frac{d}{d'} \right) \right) \\
  & = \sum_{d' \mid d \mid n} h(d') \; g\left( \frac{d}{d'} \right) \; f\left(\frac{n}{d}\right) \\
  & = \sum_{abc=n} f(a)g(b)h(c)
  \end{aligned}
  \]
  As for the identity element, consider the arithmetical function $u(n) = 1$ if and only if $n=1$. Now, $(f \star u)(n) = \sum_{d \mid n} f(n/d) u(d) = f(n)$.

  We now obtain the invertibility using the following lemma

  \begin{slemma}
  An arithmetical function is invertible under the Dirichlet convolution if and only if $f(1) \neq 0$.
  \end{slemma}

  \begin{prooflem}
  Suppose $f$ is invertible, then there is a $g$ such that $f \star g = u$. Evaluating both sides at $n=1$ gives $f(1)g(1)=1$ which shows that $f(1)\neq 0$.

  Conversely, let $f(1)\neq 0$. Construct the function $g$ as $g(1) = \frac{1}{f(1)}$ and for $n > 1$, define $g$ recursively as
  \[
  g(n) = -\frac{1}{f(1)} \sum_{d \mid n, d \neq 1} f(d) g(n/d)
  \]
  It is clear that $(f\star g)(1)=1$. Let $n \geq 2$.
  \[
  \begin{aligned}
    (f \star g)(n) & = \sum_{d \mid n}f(d)g(n/d) \\
    & = f(1)g(n) + \sum_{d \mid n, d \neq 1} f(d)g(n/d) \\
    & = f(1)g(n) - f(1)g(n) \\
    & = 0
  \end{aligned}
  \]
  This shows that the $g$ constructed is indeed the inverse of $f$ since $f \star g = u$.
  \end{prooflem}
  This shows invertibility of elements of the set. Thus the set is indeed a group under the binary operation of Dirichlet convolution.
\end{proof}

\begin{proposition}
  The set of all multiplicative arithmetical functions forms a subgroup of the above described group.
\end{proposition}

\begin{proof}
  Suppose $f,g$ are multiplicative. Let $m,n$ be two coprime positive integers. We have
  \[
  (f \star g)(mn) = \sum_{d \mid mn} f(d) g(mn/d)
  \]
  Since $m$ and $n$ are coprime, every $d \mid mn$ can be expressed as $d = pq$ where $p \mid m$ and $q \mid n$ and $p$ and $q$ are coprime. Using this and the multiplicativity of $f$ and $g$, we get
  \[
  (f \star g)(mn) = \sum_{p \mid m, q \mid n} f(p) f(q) g(m/p) g(n/q)
  \]
  The multiplicativity of the convolution follows. Thus, this set is closed under the binary operation. The identity $u$ is clearly multiplicative.

  We now show that the inverse of a multiplicative arithmetical function is multiplicative. Let $f$ be an arithmetical multiplicative function with inverse given by the function $g$.

  Note that multiplicative functions $f$ satisfy $f(1)=1$.

  Recall that $g(1) = 1/f(1) = 1$ and for $n \geq 2$ we have the recursion
  \[
  g(n) = -\frac{1}{f(1)} \sum_{d \mid n, d \neq 1} f(d) g(n/d) = - \sum_{d \mid n, d \neq 1} f(d) g(n/d) 
  \]
  We want to show that $g(mn)=g(m)g(n)$ whenever $m$ and $n$ are coprime. The base case when either $m=1$ or $n=1$ is trivial. Fix $N>1$ and assume that $g$ is multiplicative on all coprime pairs $(m',n')$ with $m'n' < N$. Let $mn=N$ with $\gcd(m,n)=1$.

  We have
  \[
  g(mn) = -\sum_{d \mid mn, d \neq 1} f(d) g(mn/d) = - \sum_{\substack{p \mid m, q \mid n \\ \gcd(p,q)=1, pq \neq 1}} f(p) f(q) g(m/p) g(n/q)
  \]
  To split the $g$, we used the inductive hypothesis which worked because $m/p$ and $n/q$ were coprime and $mn/pq < mn = N$ since $pq \neq 1$. Also note that $\gcd(p,q)=1$ is forced by $\gcd(m,n)=1$.

  We thus have
  \[
  \begin{aligned}
    g(mn) & = - \left( \sum_{p \mid m, q \mid n} f(p) f(q) g(m/p) g(n/q) - f(1)g(m)f(1)g(n) \right) \\
    & = - (f\star g)(p) (f\star g)(q) + g(m)g(n)
  \end{aligned}
  \]
  But the convolutions are both zero since $f$ and $g$ are inverses. Thus, the result that $g$ is multiplicative follows.
\end{proof}

\begin{question}
  Does the group of completely multiplicative functions form a subgroup of the group of multiplicative functions?
\end{question}

\begin{answer}
  No. Consider the functions $f(n)=n$ and $g(n)=1$.
\end{answer}

\begin{question}
  Prove the following
  \[
  \frac{1}{\zeta(s)} = \sum_{n=1}^\infty \frac{\mu(n)}{n^s}
  \]
\end{question}

\begin{answer}
  From \cref{eul_prod},
  \[
  \begin{aligned}
    \frac{1}{\zeta(s)}  & = \prod_{p \text{ prime}} \left( 1 - \frac{1}{p^s} \right) \\
    & = \left( 1- \frac{1}{2^s} \right)\left( 1- \frac{1}{3^s} \right)\left( 1- \frac{1}{5^s} \right) \cdots
  \end{aligned}
  \]
  Each term will be of the form $1/n^s$. However, it is clear that only those $n$ are covered which are not divisible by squares of prime, that is, $n$ is square-free. The constant term is $1$. The result now follows from the definition of the Möbius function.
\end{answer}

\begin{question}
  Prove the following
  \[
  -\frac{\zeta'(s)}{\zeta(s)} = \sum_{n=1}^\infty \frac{\Lambda(n)}{n^s}
  \]
\end{question}

\begin{answer}
  In the proof, $p$ shall denote any and all prime numbers, for convenience. We have,
  \[
  \log{\zeta(s)} = - \sum_p \log{\left( 1-\frac{1}{p^s} \right)}
  \]
  We use the series expansion of $\log(1-x)$ for $|x|<1$ to simplify it to the following:
  \[
  \log{\zeta(s)} = - \sum_p \sum_{m=1}^\infty \frac{1}{mp^{ms}}
  \]
  Differentiating, we get:
  \[
  - \frac{\zeta'(s)}{\zeta(s)} = \sum_p \sum_{m=1}^\infty \frac{\log{p^m}}{mp^{ms}} = \sum_p \sum_{m=1}^\infty \frac{\log{p}}{p^{ms}}
  \]
  One can compare term by term and check that this is precisely the RHS given in the question. Only those $n^s$ occur for which $n$ is a prime power $n=p^k$ and the coefficient is $\log{p}$.
\end{answer}

\end{document}